% Môn: Vật lý
% Lớp: 11
% Chương: Chương III. Điện trường
% Bài: L11C3  Bài 19. Thế năng điện
% Số câu: 9
% App đọc file này trực tiếp — sửa rồi Commit trên GitHub, bấm Đồng bộ GitHub .tex

% L11C3  Bài 19. Thế năng điện

% ===== Câu 1 =====
% ID: L11C3B19-01-TL
% Mức: NB
\dangbt{Chưa có dạng}
\begin{bt}
% ID: L11C3B19-01-TL
% Mức: NB
Hạt bụi mịn ở Bài Câu 3 dịch chuyển thẳng đứng xuống dưới $10 \mathrm{cm}$ so với vị trí ban đầu sau đó lại bị các luồng không khí nâng lên trở lại vị trí cũ. Lúc này công của điện trường đều của Trái Đất trong dịch chuyển trên của hạt bụi mịn sẽ bằng: $A\cdotA = 0$,1 $\cdot$  qE. B. $A = 0$,2 $\cdot$  qE. C. $A = 0$,1. mg. D. $A = 0$. $\nguon{ly11 kntt.pdf}$
\loigiai{
Chọn D.

Vị trí đầu và cuối trùng nhau nên $A$ = 0.
}
\end{bt}

% ===== Câu 2 =====
% ID: L11C3B19-01-TN
% Mức: NB
\begin{ex}
% ID: L11C3B19-01-TN
% Mức: NB
Công của lực điện trong dịch chuyển của một điện tích q trong điện trường từ điểm $M$ đến điểm $N$ không phụ thuộc vào $\nguon{ly11 kntt.pdf}$
\choice
{\True cung đường dịch chuyển.}
{điện tích q.}
{điện trường E.}
{vị trí điểm M.}
\loigiai{
Chọn A.

Công của lực điện trong dịch chuyển của một điện tích q trong điện trường từ điểm $M$ đến điểm $N$ không phụ thuộc vào cung đường dịch chuyển mà chỉ phụ thuộc vào vị trí điểm $M$ và N.
}
\end{ex}

% ===== Câu 3 =====
% ID: L11C3B19-02-TL
% Mức: TH
\begin{ex}
% ID: L11C3B19-02-TL
% Mức: TH
Đối với điện trường của một điện tích điểm $Q$, người ta tính toán được công để dịch chuyển một $Q$ điện tích q từ vô cùng về điểm $M$ cách $Q$ một khoảng r có giá trị bằng A∞ $M = q 4\pi\varepsilon r$. Hãy tính công 0 của lực điện trong dịch chuyển của điện tích q từ vi tri $M$ cách $Q$ một khoảng $1\,\text{m}$ tới vi trí $N$ cách $Q$ một khoảng $2\,\text{m}$. $\nguon{ly11 kntt.pdf}$
\choiceTF

\loigiai{
$A_{MN} = q\dfrac{Q}{4\pi\varepsilon_{0}r_{2}} - q\dfrac{Q}{4\pi\varepsilon_{0}r_{1}} = q\dfrac{Q}{4\pi\varepsilon_{0}.2} - q\dfrac{Q}{4\pi\varepsilon_{0}.1} = - q\dfrac{Q}{8\pi\varepsilon_{0}}(J)$.
}
\end{ex}

% ===== Câu 4 =====
% ID: L11C3B19-02-TN
% Mức: NB
\begin{ex}
% ID: L11C3B19-02-TN
% Mức: NB
Trong điện trường đều của Trái Đất, chọn mặt đất là mốc thế năng điện. Một hạt bụi mịn có khối lượng m, điện tích q đang lơ lửng ở độ cao h so với mặt đất. Thế năng điện của hạt bụi mịn là: $\nguon{ly11 kntt.pdf}$
\choice
{$Wt = mgh.$}
{\True $Wt = qEh.$}
{$Wt = mEh.$}
{$Wt = qgh.$}
\loigiai{
Chọn B.

Thế năng W_(t) = qEh
}
\end{ex}

% ===== Câu 5 =====
% ID: L11C3B19-03-TL
% Mức: TH
\begin{ex}
% ID: L11C3B19-03-TL
% Mức: TH
Trong điện trường của điện tích $Q$ cố định. a) Xác định thế năng điện của một electron tại điểm $M$ cách $Q$ một khoảng $2\,\text{m}$. b) Dưới tác dụng của lực điện kéo electron từ điểm $M$ và với vận tốc ban đầu bằng 0, dịch chuyển theo đường thẳng về phía điện tích $Q \gt  0$. Tính tốc độ của electron khi còn cách điện tích $Q$ một khoảng $1\,\text{m}$. $\nguon{ly11 kntt.pdf}$
\choiceTF

\loigiai{
a) $W_{M} = qV_{M} = q\dfrac{Q}{4\pi\varepsilon_{0}r} = - 1,6\cdot 10^{-19}.\dfrac{Q}{4\pi\varepsilon_{0}.2} = - 0,2 \cdot 10^{-19} \cdot \dfrac{Q}{\pi\varepsilon_{0}}(J)$.

b) Dựa vào công thức ở bài 19.8 tính được: $v_{M} = \sqrt{\dfrac{0,4}{9,1} \cdot \dfrac{Q}{\pi\varepsilon_{0}}} \cdot 10^{6}\left( {\text{m}/\text{s}} \right)$.
}
\end{ex}

% ===== Câu 6 =====
% ID: L11C3B19-03-TN
% Mức: NB
\begin{ex}
% ID: L11C3B19-03-TN
% Mức: NB
Thế năng điện của một điện tích q đặt tại điểm $M$ trong một điện trường bất kì không phụ thuộc vào $\nguon{ly11 kntt.pdf}$
\choice
{điện tích q.}
{vị trí điểm M.}
{điện trường.}
{\True khối lượng của điện tích q.}
\loigiai{
Chọn D.

Thế năng điện của một điện tích q đặt tại điểm $M$ trong một điện trường bất kì không phụ thuộc vào khối lượng của điện tích q.
}
\end{ex}

% ===== Câu 7 =====
% ID: L11C3B19-04-TL
% Mức: TH
\begin{ex}
% ID: L11C3B19-04-TL
% Mức: TH
Một ion âm OH - có khối lượng 2,833 $\cdot10^{-26}\mathrm{kg}$ được thổi ra từ máy lọc không khí với vận tốc $10 \mathrm{m/s}$, cách mặt đất $80 \mathrm{cm}$ ở nơi có điện trường của Trái Đất bằng $120 \mathrm{V/m}$. Dưới tác dụng của lực điện, sau một thời gian, người ta quan sát thấy ion đang chuyển động với vận tốc $0,5 \mathrm{m/s}$ ở vị trí cách mặt đất $1,5\,\text{m}$. Hãy xác định công cản mà môi trường đã thực hiện trong quá trình dịch chuyển của ion nói trên. $\nguon{ly11 kntt.pdf}$
\choiceTF

\loigiai{
Cơ năng lúc ban đầu của ion âm OH^{-} bằng:

$W_{1} = qEh_{1} + \dfrac{mv_{1}^{2}}{2} = \left( {- 1,6 \cdot 10^{-19}} \right) \cdot 120 \cdot 0,8 + \dfrac{2,833 \cdot 10^{-26} \cdot 10^{2}}{2} = - 1,536 \cdot 10^{-17}\text{J}$ Cơ năng lúc sau của ion âm OH^{-} bằng: $W_{2} = qEh_{2} + \dfrac{mv_{2}^{2}}{2} = \left( {- 1,6 \cdot 10^{-19}} \right)120 \cdot 1,5 + \dfrac{2,833 \cdot 10^{-26} \cdot 0,5^{2}}{2} = - 2,880 \cdot 10^{-17}\text{J}$

Độ biến thiên cơ năng sẽ bằng với công cản của môi trường trong chuyển động của ion âm từ vị trí thứ nhất đến vị trí thứ hai: A_(can) = W₂ - W₁ =  - 1, 344 ⋅ 10^{-17}J.
}
\end{ex}

% ===== Câu 8 =====
% ID: L11C3B19-04-TN
% Mức: NB
\begin{ex}
% ID: L11C3B19-04-TN
% Mức: NB
Đặt vào hai bản kim loại phẳng song song một hiệu điện thế $U = 100\$, $\mathrm{V}$. Một hạt bụi mịn có điện tích $q = +3$,2 $\cdot$  10^{-19} $C$ lọt vào chính giữa khoảng điện trường đều giữa hai bản phẳng. Coi tốc độ hạt bụi khi bắt đầu vào điện trường đều bằng 0, bỏ qua lực cản của môi trường. Động năng của hạt bụi khi va chạm với bản nhiễm điện âm bằng: $\nguon{ly11 kntt.pdf}$
\choice
{$W_0 = 6$,4 $\cdot$  10^{-17} J.}
{$Wd = 3,2\cdot10^{-17} J.$}
{\True $W\sigma = 1,6\cdot10^{-17} J.$}
{$Wd = 0\,\mathrm{J}.$}
\loigiai{
Chọn C.

Động năng hạt bụi khi va chạm với bản âm bằng công của lực điện tác dụng lên điện tích khi nó di chuyển từ vị trí ban đầu đến bản âm

$W_{d} = A = qEd' = qE\dfrac{d}{2} = q\dfrac{U}{2} = 3,2\cdot 10^{-19}.\dfrac{100}{2} = 1,6 \cdot 10^{-17}\text{J}$
}
\end{ex}

% ===== Câu 9 =====
% ID: L11C3B19-05-TN
% Mức: NB
\begin{ex}
% ID: L11C3B19-05-TN
% Mức: NB
Một điện tích q bay vào trong một điện trường đều theo phương vuông góc với đường sức điện. Trong suốt quá trình chuyển động, thế năng điện của điện tích đó $\nguon{ly11 kntt.pdf}$
\choice
{\True luôn giảm dần.}
{luôn không đổi.}
{luôn giảm dần nếu $q \gt  0 và luôn tăng dần nếu q \lt  0$.}
{luôn giảm dần nếu $q \lt  0 và luôn tăng dần nếu q \gt  0$.}
\loigiai{
Chọn A.

Thế năng của điện tích giảm dần.
}
\end{ex}
